% !TEX root = task-6-6-1.tex
% Editable task bank. Append new BankTask blocks at the end of the document.
\documentclass[12pt,a4paper]{article}
\usepackage{physics-textbook}
\usepackage{graphicx,needspace}
\usetikzlibrary{decorations.pathmorphing}
\newif\ifShowSolutions
\ifdefined\TutorEdition
  \ShowSolutionstrue
\else
  \ShowSolutionsfalse % Change to \ShowSolutionstrue to display solutions.
\fi
\newcounter{banktask}
\newsavebox{\BankTaskBox}
\newcommand{\BankTaskContent}[3]{%
  \noindent\textbf{\thebanktask.}\quad #2\par
  \ifShowSolutions
    \Needspace{50mm}% Keep the solution heading with its opening text or diagram.
    #3\par
    \ifdefined\TutorEdition
      \ifcsname TutorSource#1\endcsname
        \nopagebreak\smallskip{\small\csname TutorSource#1\endcsname\par}%
      \fi
    \fi
  \fi
}
\newcommand{\BankTask}[3]{%
  \refstepcounter{banktask}\par\medskip
  \sbox{\BankTaskBox}{\begin{minipage}{\linewidth}%
    \BankTaskContent{#1}{#2}{#3}\end{minipage}}%
  \ifdim\dimexpr\ht\BankTaskBox+\dp\BankTaskBox\relax>\dimexpr\textheight-20mm\relax
    \Needspace{35mm}\BankTaskContent{#1}{#2}{#3}%
  \else
    \Needspace{\dimexpr\ht\BankTaskBox+\dp\BankTaskBox\relax}%
    \noindent\usebox{\BankTaskBox}\par
  \fi
  \medskip
}
\begin{document}
\BookSubsection{6.6.1}{Постулаты Бора. Энергетические уровни и спектр атома водорода}


% Source sdamgia-3161: https://phys-ege.sdamgia.ru/problem?id=3161
\BankTask{sdamgia-3161}{
\TaskMatchingColumns{Постулаты Бора}{А) первый\par Б) второй}{Формулировки}{1) переходя из одного состояния в другое, атом излучает (поглощает) половину разности энергий в начальном и конечном состояниях\par 2) переходя из одного состояния в другое, атом излучает (поглощает) квант энергии, равный разности энергий в начальном и конечном состояниях\par 3) атом может находиться только в одном из двух возможных состояний\par 4) атом может находиться только в одном из состояний с определённым значением энергии}\begin{center}\begin{tabular}{|c|c|}\hline А&Б\\\hline &\\\hline\end{tabular}\end{center}
}{
\smallskip\noindent\textbf{Решение}\par
Первый постулат Бора гласит: «Атомная система может находиться только в особых стационарных, или квантовых, состояниях, каждому из которых соответствует определенная энергия» (А  — $4$). Второй постулат Бора: «Излучение (поглощение) света происходит только при переходе из одного стационарного состояния в другое, при этом излучается (поглощается) квант энергии равный разности энергий в начальном и конечном состояниях» (Б  — $2$).

Ответ: $42$.
}

% Source sdamgia-9513: https://phys-ege.sdamgia.ru/problem?id=9513
\BankTask{sdamgia-9513}{
Электрон в атоме водорода перешел с низкой стационарной орбиты (с меньшим номером $n$) на более высокую стационарную орбиту (с бóльшим номером $n$). Как в результате этого изменились модуль силы электрического взаимодействия электрона с ядром и полная энергия электрона?\par1. Увеличится. 2. Уменьшится. 3. Не изменится.\par\TaskChangeTable{Модуль силы электрического взаимодействия электрона с ядром}{Полная энергия электрона}
}{
\smallskip\noindent\textbf{Решение}\par
Сила электрического взаимодействия (сила Кулона) обратно пропорциональна квадрату расстояния между заряженными частицами. Поскольку электрон переходит на более высокую стационарную орбиту, увеличивается расстояние между электроном и положительно заряженным ядром, а следовательно, уменьшается модуль силы электрического взаимодействия. 
Уровни энергии электрона в атоме водорода задаются формулой $E_n=-13{,}6\,\text{эВ}/n^2$, где $n=1,2,3,\ldots$ При переходе электрона на более высокую орбиту его полная энергия увеличивается.

Ответ: $21$.
}

% Source sdamgia-10326: https://phys-ege.sdamgia.ru/problem?id=10326
\BankTask{sdamgia-10326}{
\TaskMatchingColumns{Физическая величина}{А) Энергия электрона, находящегося в атоме водорода на энергетическом уровне с номером $n$.\par Б) Энергия, которую нужно сообщить электрону в атоме водорода для того, чтобы он перешёл с $n$-го энергетического уровня на $m$-й энергетический уровень.}{Формула}{1) $E_{nm}=13{,}6\,\text{эВ}\left(\frac1{n^2}-\frac1{m^2}\right)$\par 2) $E_{nm}=13{,}6\,\text{эВ}\left(\frac1n-\frac1m\right)$\par 3) $E_n=-\frac{13{,}6\,\text{эВ}}{n^2}$\par 4) $E_n=-\frac{13{,}6\,\text{эВ}}n$}\begin{center}\begin{tabular}{|c|c|}\hline А&Б\\\hline &\\\hline\end{tabular}\end{center}
}{
\smallskip\noindent\textbf{Решение}\par
Уровни энергии электрона в атоме водорода даются формулой 

\[E_n= \frac{- 13{,}6\,\text{эВ}}{n ^2} ,\]

где $n=1,2,3,\ldots$ Чтобы атом перешел из состояния $E_{n}$ в состояние $E_{m}$ ему надо сообщить энергию 

\[E_{nm}=E_m - E_n=13{,}6\,\text{эВ} ( \frac{1}{n ^2} - \frac{1}{m ^2} ) .\]

Ответ: $31$.
}

% Source sdamgia-6204: https://phys-ege.sdamgia.ru/problem?id=6204
\BankTask{sdamgia-6204}{
В таблице приведены значения энергии для второго и четвертого энергетических уровней атома водорода.

\begin{center}\begin{tabular}{|c|c|}\hline
Номер уровня&Энергия, $10^{-19}$ Дж\\\hline
$2$&$-5{,}45$\\\hline
$4$&$-1{,}36$\\\hline
\end{tabular}\end{center}

Какой должна быть энергия фотона, при поглощении которого атом переходит со второго уровня на четвертый? (в $10^{-19}$ Дж.)
}{
\smallskip\noindent\textbf{Решение}\par
Для того, чтобы электрон перешел со второго уровня на четвертый под действием фотона, энергия фотона должна быть равна разности энергий четвертого и второго энергетических уровней: $[-1{,}36-(-5{,}45)]\cdot10^{-19}\,\text{Дж}=4{,}09\cdot10^{-19}\,\text{Дж}$.

Ответ: $4{,}09$.
}

% Source sdamgia-12869: https://phys-ege.sdamgia.ru/problem?id=12869
\BankTask{sdamgia-12869}{
Электрон в атоме водорода перешел с энергетического уровня с номером m в основное энергетическое состояние с номером n = $1$. При этом был испущен фотон с импульсом $5{,}44$·$10$$^{-27}$ кг·м/с. Чему равен номер m?
}{
\smallskip\noindent\textbf{Решение}\par
Энергия испущенного фотона равна $E=pc=5{,}44 \cdot 10 ^{ - 27 } \cdot 3 \cdot 10 ^{8} =16{,}32 \cdot 10 ^{ - 19 } \,\text{Дж}= 10{,}2\,\text{эВ}.$
Энергетические уровни водорода рассчитываются по формуле $E_n= - \frac{13{,}6\,\text{эВ}}{n ^2} .$ Получаем:

\[E_m - E_n=E \Longleftrightarrow - 13{,}6 \cdot ( \frac{1}{m ^2} - \frac{1}{1 ^2} ) =10{,}2 \Longleftrightarrow m=2.\]

Ответ: $2$.
}

% Source sdamgia-8951: https://phys-ege.sdamgia.ru/problem?id=8951
\BankTask{sdamgia-8951}{
Электрон в атоме водорода находится в основном (самом низком, с номером n = $1$) энергетическом состоянии. Атом поглощает фотон с импульсом $6{,}8\cdot10^{-27}$ кг·м/с. Найдите номер энергетического уровня, на который в результате этого перейдет электрон.
}{
\smallskip\noindent\textbf{Решение}\par
Уровни энергии электрона в атоме водорода задаются формулой $E_n=-\frac{13{,}6\,\text{эВ}}{n^2}$, где $n=1,2,3,\ldots$ При поглощении фотона атом переходит из основного состояния $E_1$ в некоторое состояние $E_n.$ Зная импульс фотона, можем найти его энергию:

\[E=p \cdot c ( \,\text{Дж} ) = \frac{p \cdot c}{1{,}6 \cdot 10 ^{ - 19 }} ( \,\text{эВ} ) ,\]

где $c$ — скорость света. Запишем закон сохранения энергии:

\[\frac{p \cdot c}{1{,}6 \cdot 10 ^{ - 19 }} =E_n - E_1=13{,}6 \cdot ( 1 - \frac{1}{n ^2} ) \Rightarrow n ^2 = \frac{1}{1 - \frac{p \cdot c}{1{,}6 \cdot 10 ^{ - 19 } \cdot 13{,}6}} =16.\]

Отсюда находим, что электрон перейдет на уровень с номером $n=4$.

Ответ: $4$.
}

% Source sdamgia-10650: https://phys-ege.sdamgia.ru/problem?id=10650
\BankTask{sdamgia-10650}{
Электрон в атоме водорода переходит с энергетического уровня с номером n = $2$ на энергетический уровень с n = $1$. Чему равен модуль импульса испущенного при этом фотона? (В $10^{-29}\,\text{кг}\cdot\text{м/с}$; округлите до целого числа.)
}{
\smallskip\noindent\textbf{Решение}\par
Уровни энергии электрона в атоме водорода задаются формулой $E_n=-13{,}6\,\text{эВ}/n^2$, где $n=1,2,3,\ldots$ При переходе атома из состояния $E_{n}$ в состояние $E_{m}$ атом испускает фотон с энергией $E=E_n - E_m.$ При переходе со второго на первый энергетический уровень атом испустит фотон с энергией

\[E=E_2 - E_1=13{,}6\,\text{эВ}\cdot(1-\frac14) =10{,}2\,\text{эВ}.\]

Импульс фотона связан с его энергией как

\[p=\frac{E}{c}\Longleftrightarrow 10^{29}\cdot p=10^{29}\cdot\frac{10{,}2\cdot1{,}6\cdot10^{-19}}{3\cdot10^8}=544\,\text{кг}\cdot\text{м/с}.\]

Ответ: $544$.
}

% Source sdamgia-4898: https://phys-ege.sdamgia.ru/problem?id=4898
\BankTask{sdamgia-4898}{
Уровни   энергии   электрона   в   атоме    водорода   задаются    формулой

$E= - \frac{13{,}6}{n ^2}$ эВ, где $n = 1, 2, 3, ...$ . При переходе атома из состояния $E_2$ в состояние $E_1$ атом испускает фотон. Попав на поверхность фотокатода, этот фотон выбивает фотоэлектрон. Частота света, соответствующая красной границе фотоэффекта для материала поверхности фотокатода, $\nu_{\text{кр}} = 6 \cdot 10^{14}$ Гц. Чему равна максимальная возможная кинетическая энергия фотоэлектрона?
}{
\smallskip\noindent\textbf{Решение}\par
Согласно постулатам Бора, свет излучается при переходе атома на более низкие уровни энергии, при этом фотоны несут энергию, равную разности энергий начального и конечного состояний. Таким образом, испущенный фотон имел нес энергию 

\[h\nu=E_2 - E_1= - 13{,}6 \,\text{эВ} ( \frac{1}{2 ^2} - \frac{1}{1 ^2} ) = 10{,}2 \,\text{эВ}\]

Согласно уравнению фотоэффекта, максимальная кинетическая энергия вылетающих фотоэлектронов связана с энергией фотона и работой выхода соотношением

\[h\nu=A + E\]

Работа выхода связана с частотой красной границы соотношением: 

\[A=h\nu_{\,\text{кр}}\]

Таким образом, максимально возможная кинетическая энергия фотоэлектрон равна

\[\begin{aligned}
E&=h\nu-A=10{,}2\,\text{эВ}-6{,}62\cdot 10^{-34}(\text{Дж}\cdot\text{с})\\
 &\quad\cdot 6\cdot10^{14}\,\text{Гц}\cdot\frac{1}{1{,}6}\cdot 10^{19}\,\text{эВ}\approx 7{,}7\,\text{эВ}.
\end{aligned}\]

Ответ: $\approx 7{,}7 \,\text{эВ}$
}

% Source sdamgia-3811: https://phys-ege.sdamgia.ru/problem?id=3811
\BankTask{sdamgia-3811}{

\begin{center}\begin{tikzpicture}[book diagram,>=stealth]\draw[task axis,->](0,-.3)--(0,5.2)node[above]{$E$};\draw[thick](0,0)--(3.5,0);\node[left]at(0,0){$ E_0 $};\draw[thick](0,1.1)--(3.5,1.1);\node[left]at(0,1.1){$ E_1 $};\draw[thick](0,2.4)--(3.5,2.4);\node[left]at(0,2.4){$ E_2 $};\draw[thick](0,3.4)--(3.5,3.4);\node[left]at(0,3.4){$ E_3 $};\draw[thick](0,4.1)--(3.5,4.1);\node[left]at(0,4.1){$ E_4 $};\draw[->]( 0.7,0)--(0.7,3.4)node[midway,right]{$ 1 $};\draw[->]( 1.4,0)--(1.4,4.1)node[midway,right]{$ 2 $};\draw[->]( 2.1,2.4)--(2.1,0)node[midway,right]{$ 3 $};\draw[->]( 2.8,1.1)--(2.8,0)node[midway,right]{$ 4 $};\draw[dashed](0,5)--(3.5,5);\node[left]at(0,5){$0$};\end{tikzpicture}\end{center}
На рисунке изображена упрощённая диаграмма энергетических уровней атома. Нумерованными стрелками отмечены некоторые возможные переходы атома между этими уровнями. \TaskMatchingColumns{Процесс}{А) Поглощение света наибольшей длины волны\par Б) Испускание света наибольшей длины волны}{Стрелка}{1) 1\par 2) 2\par 3) 3\par 4) 4}\begin{center}\begin{tabular}{|c|c|}\hline А&Б\\\hline &\\\hline\end{tabular}\end{center}
}{
\smallskip\noindent\textbf{Решение}\par
Квант энергии поглощается при переходе на более высокий уровень, излучается при переходе на более низкий. При этом расстояние между уровнями обратно пропорционально длине волны: $\Delta E= \frac{hc}{\lambda} .$ Таким образом, поглощению света наибольшей длины волны соответствует энергетический переход $1$, а излучению света наибольшей длины волны отвечает переход $4$.

Ответ: $14$.
}

% Source sdamgia-4495: https://phys-ege.sdamgia.ru/problem?id=4495
\BankTask{sdamgia-4495}{

\begin{center}\begin{tikzpicture}[book diagram,>=stealth]\draw[task axis,->](0,-.3)--(0,5.2)node[above]{$E,\,\text{Дж}$};\draw[thick](0,4.6)--(3.5,4.6);\node[left]at(0,4.6){$ 0 $};\draw[thick](0,3.6)--(3.5,3.6);\node[left]at(0,3.6){$ -2\cdot10^{-19} $};\draw[thick](0,2.2)--(3.5,2.2);\node[left]at(0,2.2){$ -5\cdot10^{-19} $};\draw[thick](0,0.5)--(3.5,0.5);\node[left]at(0,0.5){$ -8\cdot10^{-19} $};\node[right]at(3.5,3.6){$E^{(3)}$};\node[right]at(3.5,2.2){$E^{(2)}$};\node[right]at(3.5,.5){$E^{(1)}$};\end{tikzpicture}\end{center}
На рисунке изображена схема низших энергетических уровней атома. В начальный момент времени атом находится в состоянии с энергией $E^{(2)}$. Согласно постулатам Бора с какой энергией данный атом может излучать фотоны? (В $10^{-19}\,\text{Дж}$.)
}{
\smallskip\noindent\textbf{Решение}\par
Согласно постулатам Бора, свет излучается при переходе атома на более низкие уровни энергии, при этом фотоны несут энергию, равную разности энергий начального и конечного состояний. Из приведенной здесь схемы видно, что фотон может излучиться только при переходе атома в состояние $E^{(1)} ,$ при этом его энергия будет равна $- 5 \cdot 10 ^{ - 19 } \,\text{Дж} - ( - 8 \cdot 10 ^{ - 19 } \,\text{Дж} ) =3 \cdot 10 ^{ - 19 } \,\text{Дж}.$

Ответ: $3$.
}

% Source sdamgia-31893: https://phys-ege.sdamgia.ru/problem?id=31893
\BankTask{sdamgia-31893}{
\begin{center}\begin{tikzpicture}[book diagram,>=stealth,y=1.3cm]\draw[task axis,->](0,-.3)--(0,5.9)node[above]{$E$};\draw[thick](0,0)--(3.5,0);\node[left]at(0,0){$ -E_0 $};\draw[thick](0,2.6)--(3.5,2.6);\node[left]at(0,2.6){$ -\frac{E_0}{4} $};\draw[thick](0,3.7)--(3.5,3.7);\node[left]at(0,3.7){$ -\frac{E_0}{9} $};\draw[thick](0,4.3)--(3.5,4.3);\node[left]at(0,4.3){$ -\frac{E_0}{16} $};\draw[thick](0,4.8)--(3.5,4.8);\node[left]at(0,4.8){$ -\frac{E_0}{25} $};\draw[->]( 0.7,4.8)--(0.7,0)node[midway,right]{$  $};\draw[->]( 1.5,2.6)--(1.5,4.8)node[midway,right]{$  $};\draw[->]( 2.3,3.7)--(2.3,4.3)node[midway,right]{$  $};\draw[->]( 3.1,3.7)--(3.1,2.6)node[midway,right]{$  $};\draw[dashed](0,5.4)--(3.5,5.4);\node[left]at(0,5.4){$0$};\end{tikzpicture}\end{center}
На рисунке изображена упрощенная диаграмма нижних энергетических уровней атома водорода ($E_0$  =  $13{,}6\,\text{эВ}$). Стрелками отмечены некоторые возможные переходы атома между этими уровнями.
\TaskMatchingColumns{Процесс}{А) Поглощение фотона наибольшей частоты\par Б) Излучение фотона наименьшей длины волны}{Энергия фотона}{1) $0{,}84E_0$\par 2) $0{,}21E_0$\par 3) $0{,}48E_0$\par 4) $0{,}96E_0$}\begin{center}\begin{tabular}{|c|c|}\hline А&Б\\\hline &\\\hline\end{tabular}\end{center}
}{
\smallskip\noindent\textbf{Решение}\par
А)  Поглощение фотонов происходит при переходе с низшего энергетического уровня на высший. Большей энергии, а соответственно, и большей частоте соответствует переход по стрелке $2$. По второму постулату Бора энергия поглощенного фотона равна:

\[( 2 ) E_{\,\text{ф}}= - \frac{E_0}{25} + \frac{E_0}{4} =0{,}21E_0.\]

Б)  Излучение фотонов происходит при переходе с высшего энергетического уровня на низший. Большей энергии, а соответственно, меньшей длине волны соответствует переход по стрелке $1$. По второму постулату Бора энергия излученного фотона равна:

\[( 4 ) E_{\,\text{ф}}= - \frac{E_0}{25} + E_0=0{,}96E_0.\]

Ответ: $24$.
}

% Source sdamgia-3817: https://phys-ege.sdamgia.ru/problem?id=3817
\BankTask{sdamgia-3817}{
Уровни энергии электрона в атоме водорода задаются формулой $E_n= - \frac{13{,}6}{n ^2}$ эВ, где n = $1$, $2$, $3$, … При переходе из состояния $E_2$ в состояние $E_1$ атом испускает фотон. Поток таких фотонов падает на поверхность фотокатода. Запирающее напряжение для фотоэлектронов, вылетающих с поверхности фотокатода, $U_{\,\text{зап}}=7{,}4\,\text{В}.$ Какова работа выхода $A_{\,\text{вых}}$ фотоэлектронов с поверхности фотокатода?
}{
\smallskip\noindent\textbf{Решение}\par
Энергия фотона равна $h\nu = E_2 - E_1.$ Используя уравнение Эйнштейна для фотоэффекта: $h\nu = A_{\,\text{вых}} + eU_{\,\text{зап}},$ получаем $A_{\,\text{вых}} = ( E_2 - E_1 ) - eU_{\,\text{зап}} = 2{,}8\,\text{эВ} \approx 4{,}5 \cdot 10 ^{ -19 } \,\text{Дж}.$

Ответ: $\,\text{А}_{\,\text{вых}} = 2{,}8\,\text{эВ} \approx 4{,}5 \cdot 10 ^{ -19 } \,\text{Дж}.$
}

% Source sdamgia-3045: https://phys-ege.sdamgia.ru/problem?id=3045
\BankTask{sdamgia-3045}{

\begin{center}\begin{tikzpicture}[book diagram,>=stealth]\draw[task axis,->](0,-.3)--(0,5.2)node[above]{$E,\,\text{эВ}$};\draw[thick](0,4.6)--(3.5,4.6);\node[left]at(0,4.6){$ 0 $};\draw[thick](0,3.6)--(3.5,3.6);\node[left]at(0,3.6){$ -2 $};\draw[thick](0,2.2)--(3.5,2.2);\node[left]at(0,2.2){$ -5 $};\draw[thick](0,0.5)--(3.5,0.5);\node[left]at(0,0.5){$ -8{,}5 $};\node[right]at(3.5,3.6){$E^{(2)}$};\node[right]at(3.5,2.2){$E^{(1)}$};\node[right]at(3.5,.5){$E^{(0)}$};\end{tikzpicture}\end{center}
Предположим, что схема нижних энергетических уровней атомов некоего элемента имеет вид, показанный на рисунке, и атомы находятся в состоянии с энергией $E^{(1)}$. Электрон, столкнувшись с одним из таких покоящихся атомов, в результате столкновения получил некоторую дополнительную энергию. Импульс электрона после столкновения с атомом оказался равным $1{,}2\cdot10^{-24}\,\text{кг}\cdot\text{м/с}$. Определите кинетическую энергию электрона до столкновения. Возможностью испускания света атомом при столкновении с электроном пренебречь. Эффектом отдачи пренебречь. (В $10^{-19}\,\text{Дж}$; округлите до десятых.)
}{
\smallskip\noindent\textbf{Решение}\par
Если при столкновении с атомом электрон приобрел энергию, то атом перешел в состояние $E^{(0)} .$ Следовательно, после столкновения кинетическая энергия электрона стала равной $E=E_0 + 3{,}5\,\text{эВ},$ где $E_0$  — энергия электрона до столкновения; отсюда: $E_0=E - 3{,}5\,\text{эВ}.$ Импульс $p$ электрона связан с его кинетической энергией соотношением $p ^2 =m ^2 v ^2 =2mE,$ или $E= \frac{p ^2}{2m} ,$ где $m$ — масса электрона.
Следовательно, $E_0= \frac{p ^2}{2m} - 3{,}5\,\text{эВ}= \frac{1{,}44 \cdot 10 ^{ - 48 }}{2 \cdot 9{,}1 \cdot 10 ^{ - 31 }} - 3{,}5 \cdot 1{,}6 \cdot 10 ^{ - 19 } \approx 2{,}3 \cdot 10 ^{ - 19 } ( \,\text{Дж} ) .$

Ответ: $E_0\approx 2{,}3 \cdot 10 ^{ - 19 } ( \,\text{Дж} ) .$
}

% Source sdamgia-3047: https://phys-ege.sdamgia.ru/problem?id=3047
\BankTask{sdamgia-3047}{
На рисунке изображены энергетические уровни атома и указаны длины волн фотонов, излучаемых и поглощаемых при переходах с одного уровня на другой. Какова длина волны фотонов, излучаемых при переходе с уровня $E_4$ на уровень $E_1$ если $\lambda_{13}=400\,\text{нм},$ $\lambda_{24}=500\,\text{нм},$ $\lambda_{32}=600\,\text{нм}$?
\begin{center}\begin{tikzpicture}[book diagram,>=stealth]\draw[task axis,->](0,-.3)--(0,5.2)node[above]{$E$};\draw[thick](0,0.3)--(3.5,0.3);\node[left]at(0,0.3){$ 1 $};\draw[thick](0,2.3)--(3.5,2.3);\node[left]at(0,2.3){$ 2 $};\draw[thick](0,3.7)--(3.5,3.7);\node[left]at(0,3.7){$ 3 $};\draw[thick](0,4.5)--(3.5,4.5);\node[left]at(0,4.5){$ 4 $};\draw[->]( 0.8,0.3)--(0.8,3.7)node[midway,right]{$ \lambda_{13} $};\draw[->]( 1.9,2.3)--(1.9,4.5)node[midway,right]{$ \lambda_{24} $};\draw[->]( 3.1,3.7)--(3.1,2.3)node[midway,right]{$ \lambda_{32} $};\end{tikzpicture}\end{center}
}{
\smallskip\noindent\textbf{Решение}\par
Частота фотона, испускаемого или поглощаемого атомом при переходе с одного уровня энергии на другой, пропорциональна разности энергий этих уровней. Поэтому имеем: $\nu_{41}=\nu_{43} + \nu_{31},$ $\nu_{43}=\nu_{42} - \nu_{32}.$
Отсюда (так как $\nu_{nk}=\nu_{kn}$): $\nu_{41}=\nu_{13} + \nu_{24} - \nu_{32}.$
Имеем:

$\nu_{13}= \frac{c}{\lambda_{13}} = \frac{3 \cdot 10 ^{ 8 }}{4 \cdot 10 ^{ - 7 }} =0{,}75 \cdot 10 ^{ 15 } \,\text{Гц}$
$\nu_{24}= \frac{c}{\lambda_{24}} = \frac{3 \cdot 10 ^{ 8 }}{5 \cdot 10 ^{ - 7 }} =0{,}6 \cdot 10 ^{ 15 } \,\text{Гц},$
$\nu_{32}= \frac{c}{\lambda_{32}} = \frac{3 \cdot 10 ^{ 8 }}{6 \cdot 10 ^{ - 7 }} =0{,}5 \cdot 10 ^{ 15 } \,\text{Гц},$
$\nu_{41}= ( 0{,}75 + 0{,}6 - 0{,}5 ) \cdot 10 ^{ 15 } \,\text{Гц}=0{,}85 \cdot 10 ^{ 15 } \,\text{Гц}.$
$\lambda_{41}= \frac{c}{\nu_{41}} = \frac{3 \cdot 10 ^{ 8 }}{0{,}85 \cdot 10 ^{ 15 }} \,\text{м}\approx 350\,\text{нм}.$

Ответ: $350\,\text{нм}$.
}

% Source sdamgia-3048: https://phys-ege.sdamgia.ru/problem?id=3048
\BankTask{sdamgia-3048}{
На рисунке представлены энергетические уровни электронной оболочки атома и указаны частоты фотонов, излучаемых и поглощаемых при переходах с одного уровня на другой. Какова длина волны фотонов, поглощаемых при переходе с уровня $E_1$ на уровень $E_4$ если $\nu_{13}=6 \cdot 10 ^{ 14 } \,\text{Гц},$ $\nu_{24}=4 \cdot 10 ^{ 14 } \,\text{Гц},$ $\nu_{32}=3 \cdot 10 ^{ 14 } \,\text{Гц}.$
\begin{center}\begin{tikzpicture}[book diagram,>=stealth]\draw[task axis,->](0,-.3)--(0,5.2)node[above]{$E$};\draw[thick](0,0.3)--(3.5,0.3);\node[left]at(0,0.3){$ 1 $};\draw[thick](0,2.3)--(3.5,2.3);\node[left]at(0,2.3){$ 2 $};\draw[thick](0,3.7)--(3.5,3.7);\node[left]at(0,3.7){$ 3 $};\draw[thick](0,4.5)--(3.5,4.5);\node[left]at(0,4.5){$ 4 $};\draw[->]( 0.8,0.3)--(0.8,3.7)node[midway,right]{$ \nu_{13} $};\draw[->]( 1.9,2.3)--(1.9,4.5)node[midway,right]{$ \nu_{24} $};\draw[->]( 3.1,3.7)--(3.1,2.3)node[midway,right]{$ \nu_{32} $};\end{tikzpicture}\end{center}
}{
\smallskip\noindent\textbf{Решение}\par
Частота фотона, испускаемого или поглощаемого атомом при переходе с одного уровня энергии на другой, пропорциональна разности энергий этих уровней: $\nu_{21}= \frac{|E_2 - E_1|}{h} .$
Поэтому запишем: $\nu _{14}=\nu _{13} + \nu _{24} - \nu _{32}= ( 6 + 4 - 3 ) \cdot 10 ^{ 14 } =7 \cdot 10 ^{ 14 } \,\text{Гц}.$ Отсюда: $\lambda _{14}= \frac{c}{\nu_{14}} .$

Ответ: $\lambda _{14}\approx 4{,}3 \cdot 10 ^{ - 7 } \,\text{м}.$
}

% Source sdamgia-4758: https://phys-ege.sdamgia.ru/problem?id=4758
\BankTask{sdamgia-4758}{
Уровни энергии электрона в атоме водорода задаются формулой $E_n= - \frac{13{,}6}{n ^2} \,\text{эВ},$ где $n = 1, 2, 3, ...$ . При переходе атома из состояния $E_2$ в состояние $E_1$ атом испускает фотон. Попав на поверхность фотокатода, фотон выбивает фотоэлектрон. Длина волны света, соответствующая красной границе фотоэффекта для материала поверхности фотокатода, 	$\lambda_{\text{кр}} = 300\,\text{нм}.$ Чему равна максимально возможная кинетическая энергия фотоэлектрона?
}{
\smallskip\noindent\textbf{Решение}\par
Согласно постулатам Бора, свет излучается при переходе атома на более низкие уровни энергии, при этом фотоны несут энергию, равную разности энергий начального и конечного состояний. Таким образом, испущенный фотон имел энергию

\[h\nu=E_2 - E_1= - 13{,}6 \,\text{эВ} ( \frac{1}{2 ^2} - \frac{1}{1 ^2} ) = 10{,}2\,\text{эВ}.\]

Согласно уравнению фотоэффекта, максимальная кинетическая энергия вылетающих фотоэлектронов связана с энергией фотона и работой выхода соотношением

\[h\nu=A + E_{\,\text{к}}.\]

Работа выхода связана с длиной волны красной границы соотношением: 

\[A= \frac{hc}{\lambda_{\,\text{кр}}}\]

Таким образом, максимально возможная кинетическая энергия фотоэлектрон равна

\[E_{\,\text{к}}=h\nu - \frac{hc}{\lambda_{\,\text{кр}}} = 10{,}2 \,\text{эВ} - \frac{4{,}1 \cdot 10 ^{ - 15 } \,\text{эВ} \cdot \,\text{с} \cdot 3 \cdot 10 ^{8} \,\text{м/с}}{3 \cdot 10 ^{ - 7 } \,\text{м}} = 6{,}1 \,\text{эВ}.\]

Ответ: $E_{\,\text{к}} = 6{,}1\,\text{эВ}$
}

% Source sdamgia-5388: https://phys-ege.sdamgia.ru/problem?id=5388
\BankTask{sdamgia-5388}{
Электрон, имеющий импульс $p=2\cdot10^{-24}\,\text{кг}\cdot\text{м/с}$ сталкивается с покоящимся протоном, образуя атом водорода в состоянии с энергией $E_{n}$ ($n = 2$). В процессе образования атома излучается фотон. Найдите частоту $\nu$ этого фотона, пренебрегая кинетической энергией атома. Уровни энергии электрона в атоме водорода задаются формулой $E_n= - \frac{13{,}6}{n ^2} \,\text{эВ},$ где $n=1,2,3,... .$
}{
\smallskip\noindent\textbf{Решение}\par
Поскольку кинетической энергией атома пренебрегаем, то вся кинетическая энергия электрона $E_0= \frac{p ^2}{2m_e}$ перейдет в энергию атома $E_2= - \frac{13{,}6}{2 ^2} \,\text{эВ}$ и энергию фотона.

Таким образом, по закону сохранения энергии энергия фотона равна 

\[E_{\,\text{фот}} =E_0 - E_2= \frac{p ^2}{2m_e} + \frac{13{,}6\,\text{эВ}}{2 ^2} =h\nu.\]

Отсюда находим частоту фотона:

\[\nu= \frac{p ^2 / ( 2m_e ) + (13{,}6\,\text{эВ})/2 ^2}{h}\]

 Тогда 

$\nu= \frac{4 \cdot 10 ^{ - 48 } \,\text{кг} \cdot \,\text{Дж}/ ( 2 \cdot 9{,}1 \cdot 10 ^{ - 31 } \,\text{кг} ) + 13{,}6/4 \cdot 1{,}6 \cdot 10 ^{ - 19 } \,\text{Дж/эВ} \cdot \,\text{эВ}}{6{,}62 \cdot 10 ^{ - 34 } \,\text{Дж} \cdot \,\text{с}} =$

$=4{,}14 \cdot 10 ^{ 15 } \,\text{Гц}\approx 4{,}1 \cdot 10 ^{ 15 } \,\text{Гц}.$

Ответ: $\nu=4{,}1 \cdot 10 ^{ 15 } \,\text{Гц}.$
}

% Source sdamgia-6513: https://phys-ege.sdamgia.ru/problem?id=6513
\BankTask{sdamgia-6513}{
Значения энергии электрона в атоме водорода задаются формулой $E_n= \frac{- 13{,}6\,\text{эВ}}{n ^2} ,n=1,2,3,....$ При переходе с верхнего уровня энергии на нижний атом излучает фотон. Переходы с верхних уровней на уровень с $n=1$ образуют серию Лаймана; на уровень с $n=2$  — серию Бальмера; на уровень с $n = 3$  — серию Пашена и так далее. Найдите отношение $\beta$ минимальной частоты фотона в серии Бальмера к максимальной частоте фотона в серии Пашена.
}{
\smallskip\noindent\textbf{Решение}\par
Частота фотона связана с его энергией равенством $h\nu=E,$ где $h$  — постоянная Планка. В серии Бальмера энергия фотона равна $E_n - E_2,$ где $n=3, 4, ....$ Аналогично в серии Пашена энергия фотона равна $E_n - E_3,$ где $n=4, 5, ....$ Минимальной частота фотона в серии Бальмера будет при условии перехода с $3$-го уровня, максимальной частота фотона в серии Пашена будет при переходе с самого высокого ($n = \infty$) уровня. Поэтому 

\[\beta = \frac{E_3 - E_2}{E_{\infty} - E_3} = \frac{\frac{1}{2 ^2} - \frac{1}{3 ^2}}{\frac{1}{3 ^2} - 0} =1{,}25.\]

Ответ: $1{,}25$.
}

% Source sdamgia-11814: https://phys-ege.sdamgia.ru/problem?id=11814
\BankTask{sdamgia-11814}{
Значения энергии электрона в атоме водорода задаются формулой $E_n = - \frac{13{,}6\,\text{эВ}}{n ^2} ,$ где $n = 1{,}2{,}3,....$ При переходах с верхних уровней энергии на нижние атом излучает фотон. Переходы с верхних уровней на уровень с n  =  $1$ образуют серию Лаймана; на уровень с n  =  $2$  — серию Бальмера; на уровень с n  =  $3$  — серию Пашена и так далее. Найдите отношение $\beta$ минимальной частоты фотона в серии Лаймана к максимальной частоте фотона в серии Бальмера.
}{
\smallskip\noindent\textbf{Решение}\par
Частота фотона тем больше, чем больше энергия перехода электрона. Частота в серии Лаймана будет минимальна при переходе с уровня ближайшего к первому, то есть со второго уровня. Энергия перехода будет равна: $E_{\text{Л,min}} = E_2 - E_1= - \frac{13{,}6\,\text{эВ}}{2 ^2} - ( - \frac{13{,}6\,\text{эВ}}{1} ) = 10{,}2\,\text{эВ}.$ Частота в серии Бальмера будет максимальна при переходе с как можно более высокого уровня на второй. Самый высокий уровень соответствует переходу из бесконечности, то есть с уровня на котором $E_{\,\text{бесконечность}} = 0.$ Тогда энергия перехода будет равна $E_{\text{Б,max}}=E_{\,\text{бесконечность}} - E_2 = 0 - ( - \frac{13{,}6\,\text{эВ}}{2 ^2} ) =3{,}4\,\text{эВ}.$ Энергия и частота фотона связаны формулой: $E = h\nu,$ следовательно, отношение частот равно отношению энергий. Получаем: $\beta = \frac{E_{\text{Л,min}}}{E_{\text{Б,max}}} = 3.$

Ответ: $3$.
}

% Source sdamgia-11575: https://phys-ege.sdamgia.ru/problem?id=11575
\BankTask{sdamgia-11575}{
В атоме водорода уровни энергии описываются формулой $E_n= - \frac{E_0}{n ^2} .$ Наибольшая длина волны излучаемого фотона при переходе на второй уровень энергии равна $655\,\text{нм}$. Найдите $E_0$.
}{
\smallskip\noindent\textbf{Решение}\par
$1$.  Структуру атомных уровень атома водорода можно схематично представить в виде следующей схемы.
\begin{center}\begin{tikzpicture}[book diagram,>=stealth]\draw[->](0,-.3)--(0,5);\draw(0,0.0)--(3.5,0.0)node[right]{$ 1 $};\draw(0,0.65)--(3.5,0.65)node[right]{$ 2 $};\draw(0,1.3)--(3.5,1.3)node[right]{$ 3 $};\draw(0,1.9500000000000002)--(3.5,1.9500000000000002)node[right]{$ 4 $};\draw[->](1,0.0)--(1,0.65);\draw[->](2,1.3)--(2,0.65);\draw(0,4.5)--(3.5,4.5);\node[left]at(.5,-.5){поглощение};\node[right]at(4.1,.95){излучение};\draw[->](.5,-.35)--(1,.3);\draw[->](4,.95)--(2.05,.975);\end{tikzpicture}\end{center}
Поглощению фотона соответствует переход с нижнего уровня на верхний. Излучению  — переход с верхнего уровня на нижний.
$2$.  В условии говорится о излучении при переходе на уровень $2$. То есть это переходы $3$→$2$, $4$→$2$, $5$→$2$. Переход с наибольшей длиной волны соответствует переходу с наименьшей частотой излучаемого фотона. А это возможно при переходе с ближайшего ко $2$ уровню энергии  — то есть с третьего на второй.$\lambda_{\max}$→$\nu_{\min},$ так как $\lambda\nu=c.$ $\nu_{\min}:$ $3$→$2$.
\begin{center}\begin{tikzpicture}[book diagram,>=stealth]\draw[->](0,-.3)--(0,5);\draw(0,0.0)--(3.5,0.0)node[right]{$ 1 $};\draw(0,0.65)--(3.5,0.65)node[right]{$ 2 $};\draw(0,1.3)--(3.5,1.3)node[right]{$ 3 $};\draw(0,1.9500000000000002)--(3.5,1.9500000000000002)node[right]{$ 4 $};\draw(0,2.6)--(3.5,2.6)node[right]{$ 5 $};\draw(0,3.25)--(3.5,3.25)node[right]{$ 6 $};\draw[->](1,1.3)--(1,0.65);\draw[->](1.5,1.9500000000000002)--(1.5,0.65);\draw[->](2,2.6)--(2,0.65);\draw[->](2.5,3.25)--(2.5,0.65);\end{tikzpicture}\end{center}
$3$.  Согласно постулатам Бора частота излучаемого фотона определяется соотношением:

\[h\nu=E_3 - E_2 \Longleftrightarrow \frac{hc}{\lambda_2} = - \frac{E_0}{3 ^2} + \frac{E_0}{2 ^2} \Longleftrightarrow \frac{hc}{\lambda_2} = \frac{9E_0 - 4E_0}{4 \cdot 9} \Longleftrightarrow \frac{hc}{\lambda_2} = \frac{5E_0}{4 \cdot 9} .\]

$4$.  Найдем значение $E_0$:

\[\frac{hc}{\lambda_2} = \frac{5E_0}{4 \cdot 9} \Longleftrightarrow E_0= \frac{36hc}{5\lambda_2} = \frac{36 \cdot 6{,}6 \cdot 10 ^{ - 34 } \cdot 3 \cdot 10 ^{8}}{5 \cdot 655 \cdot 10 ^{ - 9 }} =2{,}18 \cdot 10 ^{ - 18 } \,\text{Дж}.\]

Ответ: $2{,}18 \cdot 10 ^{ - 18 } \,\text{Дж}.$
}


% Source F926FD: https://ege.fipi.ru/bank/questions.php?search=1&proj=BA1F39653304A5B041B656915DC36B38&qid=F926FD
\BankTask{F926FD}{
На рисунке изображена упрощённая диаграмма нижних энергетических уровней атома. Стрелками отмечены некоторые возможные переходы атома между этими уровнями.

Установите соответствие между процессами поглощения света наименьшей длины волны и излучения света наименьшей частоты и энергией соответствующего фотона.

\TaskMatchingColumns[.53\linewidth]{Процесс}{
А) поглощение света наименьшей длины волны\par
Б) излучение света наименьшей частоты
}{Энергия фотона}{
1) $E_1-E_0$\par
2) $E_2-E_0$\par
3) $E_3-E_0$\par
4) $E_4-E_0$
}

\begin{tabular}{|c|c|}
\hline
А & Б \\
\hline
\phantom{0} & \phantom{0} \\
\hline
\end{tabular}

\begin{center}
\begin{tikzpicture}[book diagram,x=0.035cm,y=0.035cm]
\draw (27,0)--(27,175);
\draw[dash pattern=on 24pt off 24pt] (32,159)--(144,159);
\node[left] at (23,159) {$0$};
\draw (27,137)--(144,137);
\draw (27,121)--(144,121);
\draw (27,101)--(144,101);
\draw (27,67)--(144,67);
\draw (27,17)--(144,17);
\node[left] at (23,137) {$E_4$};
\node[left] at (23,121) {$E_3$};
\node[left] at (23,101) {$E_2$};
\node[left] at (23,67) {$E_1$};
\node[left] at (23,17) {$E_0$};
\draw[-{book tip}] (52,17)--(52,137);
\draw[-{book tip}] (75,17)--(75,101);
\draw[-{book tip}] (97,67)--(97,17);
\draw[-{book tip}] (118,121)--(118,17);
\end{tikzpicture}
\end{center}
}{
\smallskip\noindent\textbf{Решение от GPT}\par
Энергия фотона равна $\varepsilon=h\nu=hc/\lambda=\lvert E_f-E_i\rvert$. Среди показанных переходов поглощения $E_0\to E_4$ и $E_0\to E_2$ наименьшей длине волны соответствует наибольшая энергия $E_4-E_0$ (вариант 4, а не 2). Среди переходов излучения $E_1\to E_0$ и $E_3\to E_0$ наименьшей частоте соответствует меньшая энергия $E_1-E_0$ (вариант 1, а не 3). Ответ: А — 4, Б — 1.
}

% Source AB5BFF: https://ege.fipi.ru/bank/questions.php?search=1&proj=BA1F39653304A5B041B656915DC36B38&qid=AB5BFF
\BankTask{AB5BFF}{
На рисунке изображена упрощённая диаграмма нижних энергетических уровней атома. Нумерованными стрелками отмечены некоторые возможные переходы атома между этими уровнями. Какой из этих четырёх переходов связан с поглощением света наименьшей частоты, а какой — с излучением света наибольшей частоты?

Установите соответствие между процессами поглощения и испускания света и стрелками, указывающими энергетические переходы атома.

\begin{center}
\TaskMatchingColumns[0.58\linewidth]{Процесс}{
А) поглощение света наименьшей частоты\par
Б) излучение света наибольшей частоты
}{Энергетический переход}{
1) 1\par
2) 2\par
3) 3\par
4) 4
}

\medskip
\begin{tabular}{|c|c|}
\hline
А & Б \\
\hline
\phantom{3} & \phantom{4} \\
\hline
\end{tabular}
\end{center}

\begin{center}
\begin{tikzpicture}[book diagram,x=0.025cm,y=0.025cm]
\draw (0,-9) -- (0,203);
\draw[dashed] (0,189) -- (125,189);
\node[left] at (0,189) {$0$};
\draw (0,157) -- (125,157);
\node[left] at (0,157) {$E_4$};
\draw (0,141) -- (125,141);
\node[left] at (0,136) {$E_3$};
\draw (0,110) -- (125,110);
\node[left] at (0,105) {$E_2$};
\draw (0,63) -- (125,63);
\node[left] at (0,63) {$E_1$};
\draw (0,0) -- (125,0);
\node[left] at (0,0) {$E_0$};
\draw[-{book tip}] (26,0) -- (26,141);
\draw[-{book tip}] (51,110) -- (51,0);
\draw[-{book tip}] (79,0) -- (79,63);
\draw[-{book tip}] (105,157) -- (105,0);
\node[below] at (26,0) {$1$};
\node[below] at (51,0) {$2$};
\node[below] at (79,0) {$3$};
\node[below] at (105,0) {$4$};
\end{tikzpicture}
\end{center}
}{
\smallskip\noindent\textbf{Решение от GPT}\par
Частота света определяется разностью энергий уровней: $\nu=|E_f-E_i|/h$. Поглощению соответствуют переходы 1 ($E_0\to E_3$) и 3 ($E_0\to E_1$), поэтому наименьшая частота соответствует переходу 3. Излучению соответствуют переходы 2 ($E_2\to E_0$) и 4 ($E_4\to E_0$), поэтому наибольшая частота соответствует переходу 4. Ответ: А — 3, Б — 4.
}

% Source 6BBB75: https://ege.fipi.ru/bank/questions.php?search=1&proj=BA1F39653304A5B041B656915DC36B38&qid=6BBB75
\BankTask{6BBB75}{
На рисунке изображена упрощённая диаграмма нижних энергетических уровней атома. Нумерованными стрелками отмечены некоторые возможные переходы атома между этими уровнями. Какие из этих четырёх переходов связаны с излучением света с наибольшей длиной волны и поглощением света с наименьшей энергией?

Установите соответствие между процессами поглощения и излучения света и энергетическими переходами атома, указанными стрелками.

\begin{center}
\TaskMatchingColumns[0.58\linewidth]{Процессы}{
А) излучение света с наибольшей длиной волны\par
Б) поглощение света с наименьшей энергией
}{Энергетические переходы}{
1) 1\par
2) 2\par
3) 3\par
4) 4
}

\medskip
\begin{tabular}{|c|c|}
\hline
А & Б \\
\hline
\phantom{1} & \phantom{3} \\
\hline
\end{tabular}
\end{center}

\begin{center}
\begin{tikzpicture}[book diagram,x=0.025cm,y=0.025cm]
\draw (0,-13) -- (0,192);
\draw[dashed] (0,177) -- (117,177);
\node[left] at (0,177) {$0$};
\draw (0,148) -- (117,148);
\node[left] at (0,148) {$E_4$};
\draw (0,133) -- (117,133);
\node[left] at (0,128) {$E_3$};
\draw (0,104) -- (117,104);
\node[left] at (0,99) {$E_2$};
\draw (0,59) -- (117,59);
\node[left] at (0,59) {$E_1$};
\draw (0,0) -- (117,0);
\node[left] at (0,0) {$E_0$};
\draw[-{book tip}] (25,59) -- (25,0);
\draw[-{book tip}] (49,104) -- (49,0);
\draw[-{book tip}] (73,0) -- (73,133);
\draw[-{book tip}] (97,0) -- (97,148);
\node[below] at (25,0) {$1$};
\node[below] at (49,0) {$2$};
\node[below] at (73,0) {$3$};
\node[below] at (97,0) {$4$};
\end{tikzpicture}
\end{center}
}{
\smallskip\noindent\textbf{Решение от GPT}\par
Энергия фотона равна $\varepsilon=|E_f-E_i|$, а длина волны $\lambda=hc/\varepsilon$. Излучению соответствуют переходы 1 ($E_1\to E_0$) и 2 ($E_2\to E_0$); меньшая энергия и наибольшая длина волны соответствуют переходу 1. Поглощению соответствуют переходы 3 ($E_0\to E_3$) и 4 ($E_0\to E_4$); наименьшая энергия соответствует переходу 3. Ответ: А — 1, Б — 3.
}

% Source BE27DA: https://ege.fipi.ru/bank/questions.php?search=1&proj=BA1F39653304A5B041B656915DC36B38&qid=BE27DA
\BankTask{BE27DA}{
На рисунке изображена упрощённая диаграмма нижних энергетических уровней атома. Нумерованными стрелками отмечены некоторые возможные переходы атома между этими уровнями. Какой из этих четырёх переходов связан с поглощением кванта света наибольшей частоты, а какой — с излучением кванта света наименьшей частоты?

Установите соответствие между процессами поглощения и излучения света и энергетическими переходами атома, указанными стрелками.

\begin{center}
\TaskMatchingColumns[0.58\linewidth]{Процессы}{
А) поглощение кванта света наибольшей частоты\par
Б) излучение кванта света наименьшей частоты
}{Энергетические переходы}{
1) 1\par
2) 2\par
3) 3\par
4) 4
}

\medskip
\begin{tabular}{|c|c|}
\hline
А & Б \\
\hline
\phantom{4} & \phantom{1} \\
\hline
\end{tabular}
\end{center}

\begin{center}
\begin{tikzpicture}[book diagram,x=0.025cm,y=0.025cm]
\draw (0,-13) -- (0,184);
\draw[dashed] (0,170) -- (112,170);
\node[left] at (0,170) {$0$};
\draw (0,142) -- (112,142);
\node[left] at (0,142) {$E_4$};
\draw (0,127) -- (112,127);
\node[left] at (0,122) {$E_3$};
\draw (0,99) -- (112,99);
\node[left] at (0,94) {$E_2$};
\draw (0,57) -- (112,57);
\node[left] at (0,57) {$E_1$};
\draw (0,0) -- (112,0);
\node[left] at (0,0) {$E_0$};
\draw[-{book tip}] (23,57) -- (23,0);
\draw[-{book tip}] (46,99) -- (46,0);
\draw[-{book tip}] (71,0) -- (71,127);
\draw[-{book tip}] (93,0) -- (93,142);
\node[below] at (23,0) {$1$};
\node[below] at (46,0) {$2$};
\node[below] at (71,0) {$3$};
\node[below] at (93,0) {$4$};
\end{tikzpicture}
\end{center}
}{
\smallskip\noindent\textbf{Решение от GPT}\par
Частота кванта света равна $\nu=|E_f-E_i|/h$. Поглощению соответствуют переходы 3 ($E_0\to E_3$) и 4 ($E_0\to E_4$), поэтому наибольшая частота соответствует переходу 4. Излучению соответствуют переходы 1 ($E_1\to E_0$) и 2 ($E_2\to E_0$), поэтому наименьшая частота соответствует переходу 1. Ответ: А — 4, Б — 1.
}

% Source 4Ae2Ae: https://ege.fipi.ru/bank/questions.php?search=1&proj=BA1F39653304A5B041B656915DC36B38&qid=4Ae2Ae
\BankTask{4Ae2Ae}{
На рисунке изображена упрощённая диаграмма нижних энергетических уровней атома. Стрелками отмечены некоторые возможные переходы атома между этими уровнями.

Установите соответствие между процессами поглощения света наибольшей длины волны и излучения света с наименьшей энергией и энергией соответствующего фотона.

\begin{center}
\TaskMatchingColumns[0.58\linewidth]{Процесс}{
А) поглощение света наибольшей длины волны\par
Б) излучение света с наименьшей энергией
}{Энергия фотона}{
1) $E_1-E_0$\par
2) $E_2-E_0$\par
3) $E_3-E_0$\par
4) $E_4-E_0$
}
\end{center}

\begin{center}
\begin{tabular}{|c|c|}
\hline
А & Б \\
\hline
\phantom{2} & \phantom{1} \\
\hline
\end{tabular}
\end{center}

\begin{center}
\begin{tikzpicture}[book diagram, x=1cm, y=1cm]
\draw[thin] (0,-0.4) -- (0,2.55);
\draw[thin] (-0.4,0) -- (1.8,0);
\draw[dashed] (0,2.42) -- (1.8,2.42);
\draw (0,2) -- (1.55,2);
\draw (0,1.72) -- (1.55,1.72);
\draw (0,1.23) -- (1.55,1.23);
\draw (0,0.70) -- (1.55,0.70);
\node[left] at (-0.05,2.34) {$0$};
\node[left] at (-0.05,2) {$E_4$};
\node[left] at (-0.05,1.72) {$E_3$};
\node[left] at (-0.05,1.23) {$E_2$};
\node[left] at (-0.05,0.70) {$E_1$};
\node[below left] at (-0.05,0) {$E_0$};
\draw[-{book tip}] (0.40,0.70) -- (0.40,0);
\draw[-{book tip}] (0.74,2) -- (0.74,0);
\draw[-{book tip}] (0.98,0) -- (0.98,1.23);
\draw[-{book tip}] (1.30,0) -- (1.30,1.72);
\end{tikzpicture}
\end{center}
}{
\smallskip\noindent\textbf{Решение от GPT}\par
Энергия фотона равна разности энергий уровней: $\varepsilon=hc/\lambda$. Среди показанных переходов поглощения $E_0\to E_2$ и $E_0\to E_3$ наибольшей длине волны соответствует меньшая энергия $E_2-E_0$ — вариант 2. Среди переходов излучения $E_1\to E_0$ и $E_4\to E_0$ наименьшая энергия равна $E_1-E_0$ — вариант 1; варианты 3 и 4 соответствуют более высоким энергиям. Ответ: А — 2, Б — 1.
}

% Source Be82c1: https://ege.fipi.ru/bank/questions.php?search=1&proj=BA1F39653304A5B041B656915DC36B38&qid=Be82c1
\BankTask{Be82c1}{
На рисунке изображена упрощённая диаграмма нижних энергетических уровней атома. Стрелками отмечены некоторые возможные переходы атома между этими уровнями.

Установите соответствие между процессами поглощения света наибольшей длины волны и испускания света наименьшей частоты и энергией соответствующего фотона.

\begin{center}
\TaskMatchingColumns[0.58\linewidth]{Процесс}{
А) поглощение света наибольшей длины волны\par
Б) испускание света наименьшей частоты
}{Энергия фотона}{
1) $E_1-E_0$\par
2) $E_2-E_0$\par
3) $E_3-E_0$\par
4) $E_4-E_0$
}
\end{center}

\begin{center}
\begin{tabular}{|c|c|}
\hline
А & Б \\
\hline
\phantom{2} & \phantom{1} \\
\hline
\end{tabular}
\end{center}

\begin{center}
\begin{tikzpicture}[book diagram, x=1cm, y=1cm]
\draw[thin] (0,-0.4) -- (0,2.74);
\draw[thin] (-0.4,0) -- (1.92,0);
\draw (0,2.10) -- (1.62,2.10);
\draw (0,1.76) -- (1.62,1.76);
\draw (0,1.34) -- (1.62,1.34);
\draw (0,0.76) -- (1.62,0.76);
\node[left] at (-0.05,2.25) {$E_4$};
\node[left] at (-0.05,1.91) {$E_3$};
\node[left] at (-0.05,1.39) {$E_2$};
\node[left] at (-0.05,0.81) {$E_1$};
\node[below left] at (-0.05,0) {$E_0$};
\draw[-{book tip}] (0.32,0.76) -- (0.32,0);
\draw[-{book tip}] (0.66,1.76) -- (0.66,0);
\draw[-{book tip}] (0.96,0) -- (0.96,1.34);
\draw[-{book tip}] (1.34,0) -- (1.34,2.10);
\end{tikzpicture}
\end{center}
}{
\smallskip\noindent\textbf{Решение от GPT}\par
Используем $\varepsilon=h\nu=hc/\lambda$. Из показанных переходов поглощения $E_0\to E_2$ и $E_0\to E_4$ наибольшая длина волны соответствует энергии $E_2-E_0$ — варианту 2, а вариант 4 даёт меньшую длину волны. Из переходов испускания $E_1\to E_0$ и $E_3\to E_0$ наименьшая частота соответствует энергии $E_1-E_0$ — варианту 1, а вариант 3 даёт большую частоту. Ответ: А — 2, Б — 1.
}

% Source 41596C: https://ege.fipi.ru/bank/questions.php?search=1&proj=BA1F39653304A5B041B656915DC36B38&qid=41596C
\BankTask{41596C}{
На рисунке изображена упрощённая диаграмма нижних энергетических уровней атома. Нумерованными стрелками отмечены некоторые возможные переходы атома между уровнями. Какие из этих четырёх переходов связаны с поглощением света с наименьшей частотой и излучением света с наименьшей энергией фотонов?

Установите соответствие между процессами поглощения и излучения света и энергетическими переходами атома, указанными стрелками.

\begin{center}
\TaskMatchingColumns[0.56\linewidth]{Процессы}{
А) поглощение света с наименьшей частотой\par
Б) излучение света с наименьшей энергией фотонов
}{Энергетические переходы}{
1) 1\par
2) 2\par
3) 3\par
4) 4
}
\end{center}

\begin{center}
\begin{tabular}{|c|c|}
\hline
А & Б \\
\hline
\phantom{1} & \phantom{3} \\
\hline
\end{tabular}
\end{center}

\begin{center}
\begin{tikzpicture}[book diagram, x=1cm, y=1cm]
\draw[thin] (0,-0.36) -- (0,3.54);
\draw (0,0) -- (2.16,0);
\draw (0,1.08) -- (2.16,1.08);
\draw (0,1.90) -- (2.16,1.90);
\draw (0,2.46) -- (2.16,2.46);
\draw (0,2.72) -- (2.16,2.72);
\draw[dashed] (0,3.26) -- (2.16,3.26);
\node[left] at (-0.05,3.26) {$0$};
\node[left] at (-0.05,2.80) {$E_4$};
\node[left] at (-0.05,2.46) {$E_3$};
\node[left] at (-0.05,1.90) {$E_2$};
\node[left] at (-0.05,1.08) {$E_1$};
\node[left] at (-0.05,0) {$E_0$};
\draw[-{book tip}] (0.44,0) -- (0.44,1.08);
\draw[-{book tip}] (0.88,0) -- (0.88,1.90);
\draw[-{book tip}] (1.34,2.46) -- (1.34,0);
\draw[-{book tip}] (1.78,2.72) -- (1.78,0);
\node[below] at (0.44,0) {$1$};
\node[below] at (0.88,0) {$2$};
\node[below] at (1.34,0) {$3$};
\node[below] at (1.78,0) {$4$};
\end{tikzpicture}
\end{center}
}{
\smallskip\noindent\textbf{Решение от GPT}\par
Переходы 1 и 2 направлены вверх и соответствуют поглощению, причём $\nu_1=(E_1-E_0)/h<\nu_2=(E_2-E_0)/h$, поэтому А — 1. Переходы 3 и 4 направлены вниз и соответствуют излучению, причём $\varepsilon_3=E_3-E_0<\varepsilon_4=E_4-E_0$, поэтому Б — 3. Ответ: А — 1, Б — 3.
}

% Source 9F8E6C: https://ege.fipi.ru/bank/questions.php?search=1&proj=BA1F39653304A5B041B656915DC36B38&qid=9F8E6C
\BankTask{9F8E6C}{
На рисунке изображена упрощённая диаграмма нижних энергетических уровней атома. Нумерованными стрелками отмечены некоторые возможные переходы атома между этими уровнями. Какие из этих переходов связаны с поглощением кванта света наибольшей длины волны и излучением кванта света с наименьшей энергией?

Установите соответствие между процессами поглощения и испускания света и стрелками, обозначающими энергетические переходы атома.

\begin{center}
\TaskMatchingColumns[0.64\linewidth]{Процессы}{
А) поглощение кванта света наибольшей длины волны\par
Б) излучение кванта света с наименьшей энергией
}{Энергетические переходы}{
1) 1\par
2) 2\par
3) 3\par
4) 4
}
\end{center}

\begin{center}
\begin{tabular}{|c|c|}
\hline
А & Б \\
\hline
\phantom{1} & \phantom{3} \\
\hline
\end{tabular}
\end{center}

\begin{center}
\begin{tikzpicture}[book diagram, x=1cm, y=1cm]
\draw (0,-0.34) -- (0,3.62);
\draw[dashed] (0,3.34) -- (2.2,3.34);
\draw (0,2.78) -- (2.2,2.78);
\draw (0,2.50) -- (2.2,2.50);
\draw (0,1.94) -- (2.2,1.94);
\draw (0,1.10) -- (2.2,1.10);
\draw (0,0) -- (2.2,0);
\node[left] at (0,3.34) {$0$};
\node[left] at (0,2.78) {$E_4$};
\node[left] at (0,2.50) {$E_3$};
\node[left] at (0,1.94) {$E_2$};
\node[left] at (0,1.10) {$E_1$};
\node[left] at (0,0) {$E_0$};
\draw[-{book tip}] (0.46,0) -- (0.46,1.10);
\draw[-{book tip}] (0.92,0) -- (0.92,1.94);
\draw[-{book tip}] (1.38,2.50) -- (1.38,0);
\draw[-{book tip}] (1.84,2.78) -- (1.84,0);
\node[below] at (0.46,0) {$1$};
\node[below] at (0.92,0) {$2$};
\node[below] at (1.38,0) {$3$};
\node[below] at (1.84,0) {$4$};
\end{tikzpicture}
\end{center}
}{
\smallskip\noindent\textbf{Решение от GPT}\par
Энергия кванта равна разности энергий уровней: $\varepsilon=|E_f-E_i|=hc/\lambda$. Поглощению соответствуют стрелки 1 и 2; поскольку $E_1-E_0<E_2-E_0$, наибольшая длина волны соответствует стрелке 1. Излучению соответствуют стрелки 3 и 4; поскольку $E_3-E_0<E_4-E_0$, наименьшую энергию имеет квант при переходе 3. Ответ: А — 1, Б — 3.
}

% Source 667358: https://ege.fipi.ru/bank/questions.php?search=1&proj=BA1F39653304A5B041B656915DC36B38&qid=667358
\BankTask{667358}{
На рисунке изображена упрощённая диаграмма энергетических уровней атома. Нумерованными стрелками отмечены некоторые возможные переходы атома между этими уровнями. Какие из этих четырёх переходов связаны с поглощением света наибольшей длины волны и излучением света наибольшей длины волны?

Установите соответствие между процессами поглощения и испускания света и соответствующими переходами между уровнями энергии.

\begin{center}
\TaskMatchingColumns[0.64\linewidth]{Процесс}{
А) поглощение света наибольшей длины волны\par
Б) излучение света наибольшей длины волны
}{Переход}{
1) $E_0\to E_1$\par
2) $E_0\to E_2$\par
3) $E_3\to E_0$\par
4) $E_4\to E_0$
}
\end{center}

\begin{center}
\begin{tabular}{|c|c|}
\hline
А & Б \\
\hline
\phantom{1} & \phantom{3} \\
\hline
\end{tabular}
\end{center}

\begin{center}
\begin{tikzpicture}[book diagram, x=1cm, y=1cm]
\draw (0,-0.34) -- (0,3.68);
\draw[dashed] (0,3.38) -- (2.2,3.38);
\draw (0,2.82) -- (2.2,2.82);
\draw (0,2.54) -- (2.2,2.54);
\draw (0,1.98) -- (2.2,1.98);
\draw (0,1.14) -- (2.2,1.14);
\draw (0,0) -- (2.2,0);
\node[left] at (0,3.38) {$0$};
\node[left] at (0,2.82) {$E_4$};
\node[left] at (0,2.54) {$E_3$};
\node[left] at (0,1.98) {$E_2$};
\node[left] at (0,1.14) {$E_1$};
\node[left] at (0,0) {$E_0$};
\draw[-{book tip}] (0.46,0) -- (0.46,1.14);
\draw[-{book tip}] (0.94,0) -- (0.94,1.98);
\draw[-{book tip}] (1.40,2.54) -- (1.40,0);
\draw[-{book tip}] (1.88,2.82) -- (1.88,0);
\node[below] at (0.46,0) {$1$};
\node[below] at (0.94,0) {$2$};
\node[below] at (1.40,0) {$3$};
\node[below] at (1.88,0) {$4$};
\end{tikzpicture}
\end{center}
}{
\smallskip\noindent\textbf{Решение от GPT}\par
Длина волны $\lambda=hc/|E_f-E_i|$ максимальна при минимальной разности энергий уровней. Переходы 1 и 2 соответствуют поглощению, причём $E_1-E_0<E_2-E_0$, поэтому для А подходит переход 1. Переходы 3 и 4 соответствуют излучению, причём $E_3-E_0<E_4-E_0$, поэтому для Б подходит переход 3. Ответ: А — 1, Б — 3.
}

% Source 9AC530: https://ege.fipi.ru/bank/questions.php?search=1&proj=BA1F39653304A5B041B656915DC36B38&qid=9AC530
\BankTask{9AC530}{
На рисунке изображена упрощённая диаграмма энергетических уровней атома. Нумерованными стрелками отмечены некоторые возможные переходы атома между этими уровнями. Какие из этих четырёх переходов связаны с поглощением света наибольшей длины волны и излучением света с наименьшей длиной волны?

Установите соответствие между процессами поглощения и испускания света и соответствующими переходами между уровнями энергии.

\begin{center}
\TaskMatchingColumns[0.64\linewidth]{Процесс}{
А) поглощение света наибольшей длины волны\par
Б) излучение света наименьшей длины волны
}{Переход}{
1) $E_0\to E_1$\par
2) $E_0\to E_2$\par
3) $E_3\to E_0$\par
4) $E_4\to E_0$
}
\end{center}

\begin{center}
\begin{tabular}{|c|c|}
\hline
А & Б \\
\hline
\phantom{1} & \phantom{4} \\
\hline
\end{tabular}
\end{center}

\begin{center}
\begin{tikzpicture}[book diagram, x=1cm, y=1cm]
\draw (0,-0.34) -- (0,3.68);
\draw[dashed] (0,3.38) -- (2.2,3.38);
\draw (0,2.82) -- (2.2,2.82);
\draw (0,2.54) -- (2.2,2.54);
\draw (0,1.98) -- (2.2,1.98);
\draw (0,1.14) -- (2.2,1.14);
\draw (0,0) -- (2.2,0);
\node[left] at (0,3.38) {$0$};
\node[left] at (0,2.82) {$E_4$};
\node[left] at (0,2.54) {$E_3$};
\node[left] at (0,1.98) {$E_2$};
\node[left] at (0,1.14) {$E_1$};
\node[left] at (0,0) {$E_0$};
\draw[-{book tip}] (0.46,0) -- (0.46,1.14);
\draw[-{book tip}] (0.94,0) -- (0.94,1.98);
\draw[-{book tip}] (1.40,2.54) -- (1.40,0);
\draw[-{book tip}] (1.88,2.82) -- (1.88,0);
\node[below] at (0.46,0) {$1$};
\node[below] at (0.94,0) {$2$};
\node[below] at (1.40,0) {$3$};
\node[below] at (1.88,0) {$4$};
\end{tikzpicture}
\end{center}
}{
\smallskip\noindent\textbf{Решение от GPT}\par
Для кванта света $\lambda=hc/|E_f-E_i|$. Поглощению соответствуют переходы 1 и 2; меньшая разность $E_1-E_0<E_2-E_0$ даёт наибольшую длину волны, поэтому А — 1. Излучению соответствуют переходы 3 и 4; большая разность $E_4-E_0>E_3-E_0$ даёт наименьшую длину волны, поэтому Б — 4. Ответ: А — 1, Б — 4.
}

% Source 6A3D3F: https://ege.fipi.ru/bank/questions.php?search=1&proj=BA1F39653304A5B041B656915DC36B38&qid=6A3D3F
\BankTask{6A3D3F}{
На рисунке изображена упрощённая диаграмма нижних энергетических уровней атома. Нумерованными стрелками отмечены некоторые возможные переходы атома между этими уровнями. Какие из этих переходов связаны с поглощением кванта света наибольшей частоты и излучением кванта света наименьшей частоты?

Установите соответствие между процессами поглощения и излучения света и энергетическими переходами атома, указанными стрелками.

\begin{center}
\TaskMatchingColumns[0.58\linewidth]{Процессы}{
А) поглощение кванта света наибольшей частоты\par
Б) излучение кванта света наименьшей частоты
}{Энергетические переходы}{
1) 1\par
2) 2\par
3) 3\par
4) 4
}

\medskip
\begin{tabular}{|c|c|}
\hline
А & Б \\
\hline
\phantom{2} & \phantom{3} \\
\hline
\end{tabular}
\end{center}

\begin{center}
\begin{tikzpicture}[book diagram,x=0.035cm,y=0.035cm]
\draw (0,-5)--(0,181);
\draw (0,0)--(110,0);
\draw (0,55)--(110,55);
\draw (0,97)--(110,97);
\draw (0,125)--(110,125);
\draw (0,139)--(110,139);
\draw[dash pattern=on 4pt off 4pt] (7,167)--(110,167);
\node[left] at (0,0) {$E_0$};
\node[left] at (0,55) {$E_1$};
\node[left] at (0,97) {$E_2$};
\node[left] at (0,125) {$E_3$};
\node[left] at (0,145) {$E_4$};
\node[left] at (0,167) {$0$};
\draw[-{book tip}] (23,0)--(23,55);
\draw[-{book tip}] (46,0)--(46,97);
\draw[-{book tip}] (69,125)--(69,0);
\draw[-{book tip}] (92,139)--(92,0);
\node[below] at (23,0) {$1$};
\node[below] at (46,0) {$2$};
\node[below] at (69,0) {$3$};
\node[below] at (92,0) {$4$};
\end{tikzpicture}
\end{center}
}{
\smallskip\noindent\textbf{Решение от GPT}\par
Частота кванта определяется разностью энергий уровней: $\nu=|E_f-E_i|/h$. Переходы 1 и 2 соответствуют поглощению; поскольку $E_2-E_0>E_1-E_0$, наибольшая частота соответствует переходу 2. Переходы 3 и 4 соответствуют излучению; поскольку $E_3-E_0<E_4-E_0$, наименьшая частота соответствует переходу 3. Ответ: А — 2, Б — 3.
}

% Source 8DC935: https://ege.fipi.ru/bank/questions.php?search=1&proj=BA1F39653304A5B041B656915DC36B38&qid=8DC935
\BankTask{8DC935}{
На рисунке изображена упрощённая диаграмма нижних энергетических уровней атома. Нумерованными стрелками отмечены некоторые возможные переходы атома между этими уровнями. Какой из этих четырёх переходов связан с поглощением кванта света наименьшей частоты, а какой – с излучением кванта света наибольшей частоты?

Установите соответствие между процессами поглощения и излучения света и энергетическими переходами атома, указанными стрелками.

\begin{center}
\TaskMatchingColumns[0.58\linewidth]{Процессы}{
А) поглощение кванта света наименьшей частоты\par
Б) излучение кванта света наибольшей частоты
}{Энергетические переходы}{
1) 1\par
2) 2\par
3) 3\par
4) 4
}

\medskip
\begin{tabular}{|c|c|}
\hline
А & Б \\
\hline
\phantom{3} & \phantom{2} \\
\hline
\end{tabular}
\end{center}

\begin{center}
\begin{tikzpicture}[book diagram,x=0.035cm,y=0.035cm]
\draw (0,-5)--(0,139);
\draw (0,0)--(125,0);
\draw (0,45)--(125,45);
\draw (0,74)--(125,74);
\draw (0,90)--(125,90);
\draw (0,106)--(125,106);
\draw[dash pattern=on 24pt off 24pt] (0,125)--(125,125);
\node[left] at (0,0) {$E_0$};
\node[left] at (0,45) {$E_1$};
\node[left] at (0,74) {$E_2$};
\node[left] at (0,90) {$E_3$};
\node[left] at (0,106) {$E_4$};
\node[left] at (0,125) {$0$};
\draw[-{book tip}] (19,45)--(19,0);
\draw[-{book tip}] (44,74)--(44,0);
\draw[-{book tip}] (70,0)--(70,90);
\draw[-{book tip}] (94,0)--(94,106);
\node[below] at (19,0) {$1$};
\node[below] at (44,0) {$2$};
\node[below] at (70,0) {$3$};
\node[below] at (94,0) {$4$};
\end{tikzpicture}
\end{center}
}{
\smallskip\noindent\textbf{Решение от GPT}\par
Частота кванта определяется разностью энергий уровней: $\nu=|E_f-E_i|/h$. Переходы 3 и 4 соответствуют поглощению; поскольку $E_3-E_0<E_4-E_0$, наименьшая частота соответствует переходу 3. Переходы 1 и 2 соответствуют излучению; поскольку $E_2-E_0>E_1-E_0$, наибольшая частота соответствует переходу 2. Ответ: А — 3, Б — 2.
}

% Source F37670: https://ege.fipi.ru/bank/questions.php?search=1&proj=BA1F39653304A5B041B656915DC36B38&qid=F37670
\BankTask{F37670}{
На рисунке изображена упрощённая диаграмма нижних энергетических уровней атома. Нумерованными стрелками отмечены некоторые возможные переходы атома между этими уровнями. Какой из этих четырёх переходов связан с поглощением света наименьшей частоты, а какой — с излучением света наибольшей частоты?

Установите соответствие между процессами поглощения и испускания света и стрелками, указывающими энергетические переходы атома.

\begin{center}
\TaskMatchingColumns[0.60\linewidth]{Процессы}{
А) поглощение света наименьшей частоты\par
Б) излучение света наибольшей частоты
}{Энергетические переходы}{
1) 1\par
2) 2\par
3) 3\par
4) 4
}
\end{center}

\begin{center}
\begin{tabular}{|c|c|}
\hline
А & Б \\
\hline
\phantom{3} & \phantom{2} \\
\hline
\end{tabular}
\end{center}

\begin{center}
\begin{tikzpicture}[book diagram,x=0.035cm,y=0.035cm]
\draw (0,-10) -- (0,183);
\draw[dash pattern=on 4pt off 4pt] (7,169) -- (112,169);
\node[left] at (0,169) {$0$};
\draw (0,141) -- (112,141);
\node[left] at (0,141) {$E_4$};
\draw (0,127) -- (112,127);
\node[left] at (0,127) {$E_3$};
\draw (0,98) -- (112,98);
\node[left] at (0,98) {$E_2$};
\draw (0,56) -- (112,56);
\node[left] at (0,56) {$E_1$};
\draw (0,0) -- (112,0);
\node[left] at (0,0) {$E_0$};
\draw[-{book tip}] (23,56) -- (23,0);
\draw[-{book tip}] (46,98) -- (46,0);
\draw[-{book tip}] (70,0) -- (70,127);
\draw[-{book tip}] (93,0) -- (93,141);
\node[below] at (23,0) {$1$};
\node[below] at (46,0) {$2$};
\node[below] at (70,0) {$3$};
\node[below] at (93,0) {$4$};
\end{tikzpicture}
\end{center}
}{
\smallskip\noindent\textbf{Решение от GPT}\par
Частота фотона определяется соотношением $h\nu=|E_f-E_i|$. Поглощению соответствуют стрелки 3 и 4, причём $E_3-E_0<E_4-E_0$, поэтому наименьшая частота соответствует переходу 3. Излучению соответствуют стрелки 1 и 2; поскольку $E_2-E_0>E_1-E_0$, наибольшая частота соответствует переходу 2. Ответ: А — 3, Б — 2.
}

% Source 90AF7F: https://ege.fipi.ru/bank/questions.php?search=1&proj=BA1F39653304A5B041B656915DC36B38&qid=90AF7F
\BankTask{90AF7F}{
На рисунке изображена упрощённая диаграмма нижних энергетических уровней атома. Стрелками отмечены некоторые возможные переходы атома между этими уровнями.

Установите соответствие между процессами поглощения фотона наименьшей длины волны и излучения фотона наименьшей частоты и энергией соответствующего фотона.

\begin{center}
\TaskMatchingColumns[0.60\linewidth]{Процесс}{
А) поглощение фотона наименьшей длины волны\par
Б) излучение фотона наименьшей частоты
}{Энергия фотона}{
1) $E_1-E_0$\par
2) $E_2-E_0$\par
3) $E_3-E_0$\par
4) $E_4-E_0$
}
\end{center}

\begin{center}
\begin{tabular}{|c|c|}
\hline
А & Б \\
\hline
\phantom{3} & \phantom{1} \\
\hline
\end{tabular}
\end{center}

\begin{center}
\begin{tikzpicture}[book diagram,x=0.035cm,y=0.035cm]
\draw[gray] (0,-28) -- (0,140);
\draw[densely dashed] (4,135) -- (95,135);
\node[left] at (0,135) {$0$};
\draw (2,111) -- (83,111);
\node[left] at (0,111) {$E_4$};
\draw (2,95) -- (83,95);
\node[left] at (0,95) {$E_3$};
\draw (2,69) -- (83,69);
\node[left] at (0,69) {$E_2$};
\draw (2,39) -- (83,39);
\node[left] at (0,39) {$E_1$};
\draw[gray] (-20,0) -- (99,0);
\node[below left] at (0,0) {$E_0$};
\draw[-{book tip}] (23,39) -- (23,0);
\draw[-{book tip}] (41,111) -- (41,0);
\draw[-{book tip}] (55,0) -- (55,69);
\draw[-{book tip}] (73,0) -- (73,95);
\end{tikzpicture}
\end{center}
}{
\smallskip\noindent\textbf{Решение от GPT}\par
Энергия фотона равна $\varepsilon=h\nu=hc/\lambda=|E_f-E_i|$. Из показанных переходов поглощения $E_0\to E_2$ и $E_0\to E_3$ наименьшей длине волны соответствует большая энергия $E_3-E_0$ — вариант 3. Из переходов излучения $E_1\to E_0$ и $E_4\to E_0$ наименьшей частоте соответствует меньшая энергия $E_1-E_0$ — вариант 1. Ответ: А — 3, Б — 1.
}

\end{document}
